\section{Part (b): MsPacman-v5}
\label{sec:mspacman}

\subsection{Method}

A Double DQN with a dueling head learns Ms.\ Pac-Man from the game images with the prescribed $\gamma = 0.99$ and all 9 actions. One training run takes about 6 hours, so \textbf{there is no hyperparameter ladder}. Instead, a single run (seed 0) combines the insights of part (a) with the standard Atari DQN settings.

\begin{center}
\small
\begin{tabularx}{\textwidth}{@{}p{3.1cm}L@{}}
  \toprule
  \textbf{Component} & \textbf{Choice and reason} \\
  \midrule
  Environment & \texttt{ALE/MsPacman-v5} with its defaults: frameskip 4 (one \emph{agent step} = 4 frames), sticky actions with probability 0.25, time limit 108,000 frames. \\
  State & The assignment's \texttt{preprocess\_observation} (88$\times$80 int8, divided by 128), \textbf{4 stacked frames}. A single frame shows neither the direction of movement nor the ghosts, which are not drawn in every frame (Figure~\ref{fig:b-pre}). \\
  Network & Nature CNN (32@8$\times$8/4, 64@4$\times$4/2, 64@3$\times$3/1) with a \textbf{dueling head}: separate streams of 512 units for the state value $V(s)$ and the advantages $A(s, a)$, combined as $Q = V + A - \text{mean}_a A$ (2.8M parameters). In most states of Ms.\ Pac-Man, the actions are almost equivalent (part (a) already showed tiny action gaps), so learning $V$ separately is cheap and helps. \\
  Algorithm & Double DQN target, Huber loss, gradient norm clipped at 10, Adam with a constant learning rate of $10^{-4}$ ($\epsilon_{\text{Adam}} = 1.5 \cdot 10^{-4}$), batch 32. Of the advanced techniques the assignment suggests, Huber loss and an $\epsilon$ scheduler are used; a learning-rate scheduler is not, since the slow target network already keeps the training free of divergence. \\
  Target network & Polyak update with $\tau = 0.001$ after every gradient step (time constant $\approx$ 4,000 agent steps). Part (a): the slow target network was essential. \\
  Replay buffer & \textbf{1M transitions}, stored as single int8 frames (7\,GB instead of 28\,GB for stacked states). Part (a): a large buffer prevents forgetting. Uniform sampling. \\
  Updates & 50,000 random transitions before the first update, then one gradient step every 4 agent steps. \\
  Rewards & Learning reward \textbf{clipped to $\text{sign}(r) \in \{0, 1\}$}: a dot gives 10 points, a ghost up to 1,600, so the raw scale differs by 160$\times$. Q is therefore bounded by $1/(1-\gamma) = 100$. \textbf{All reported scores are raw game scores.} \\
  Lives & A lost life is \textbf{terminal for the TD target} (no bootstrapping through a death), but the game continues. One episode is one full game with 3 lives. \\
  Exploration & $\epsilon$-greedy per agent step, linear from 1.0 to 0.1 over the first 1M steps and to 0.01 at 3M steps. \\
  Length & 6,000 episodes (the upper end of the assignment's range). \\
  \bottomrule
\end{tabularx}
\end{center}

\fig{1.0}{b_preprocessing.png}{Raw frame, preprocessed frame and the stack of 4 frames the network sees. Ms.\ Pac-Man is set to the contrast value (black); the ghosts flicker between frames.}{fig:b-pre}

\textbf{Crash recovery.} Because the run lasted a night, every 250 episodes a full checkpoint (networks, optimiser, replay buffer, all random-generator states, log) is written to a new folder and then swapped in by renaming, so a crash while writing never destroys the last good checkpoint. Every episode starts from its own seed and the GPU runs in deterministic mode, so a resumed run is \textbf{bit-identical} to an uninterrupted one (Test~10).

\subsection{Training: Required Plots (i) and (ii)}

The training took \textbf{5.8 hours for 4.46M agent steps} (17.8M frames, $\approx$ 213 agent steps per second) on a laptop GPU (GTX 1650 Max-Q).

\fig{1.0}{b_training.png}{Training: (i) maximum Q-value per episode, (ii) episode reward (game score) with the moving average of the last 100 episodes, the exploration rate and the episode length.}{fig:b-train}

\begin{itemize}
  \item \textbf{The agent learns steadily:} the moving average rises from $\approx$ 250 (random play) to $\approx$ 1,850, and the mean score of the last 500 episodes is \textbf{1,848, compared with 296 for the first 500}. The best moving average is 1,958.
  \item \textbf{The games get longer:} from $\approx$ 500 to $\approx$ 900 agent steps, so the agent survives longer, not only eats faster.
  \item \textbf{Learning follows the exploration schedule:} the curve is almost flat while $\epsilon > 0.65$ (first $\sim$750 episodes), rises fastest while $\epsilon$ decays to 0.1 (episode $\approx$ 1,750), and flattens at about 1,700--1,900 after $\epsilon$ reaches 0.01 (episode $\approx$ 4,350).
  \item \textbf{(i) Max Q stays well below the bound of 100.} It rises with the score, overshoots to 57 around episode 3,200 and settles at 40--45. There is no sign of divergence.
\end{itemize}

\subsection{Model Selection and Required Plot (iii)}

As in part (a), the reported network is chosen on \textbf{validation games} (seeds from 10,000 on), separate from the 500 test games (seeds 0--499). Every snapshot (every 500 episodes) plays 30 greedy validation games, and the best one is evaluated on the test games.

\fig{0.75}{b_snapshots.png}{Greedy validation score (30 games) of the snapshots during training. The star marks the selected snapshot.}{fig:b-snap}

\fig{0.75}{b_histogram.png}{(iii) 500 greedy test games of the selected network (episode 4,000), with the random policy as reference.}{fig:b-hist}

\textbf{(iii) Evaluation of the selected network (episode 4,000): mean = 1,790.6, std = 286.8} (Figure~\ref{fig:b-hist}).

\begin{center}
\small
\begin{tabular}{@{}lrrrr@{}}
  \toprule
  \thead[l]{Network (500 greedy test games)} & \thead[r]{Mean $\pm$ std} & \thead[r]{Median} & \thead[r]{Max} & \thead[r]{Mean length} \\
  \midrule
  \textbf{Selected snapshot (episode 4,000)} & $\mathbf{1{,}790.6 \pm 286.8}$ & 1,780 & 3,300 & 924 \\
  Final weights (episode 6,000) & $1{,}842.7 \pm 624.7$ & 1,715 & 4,740 & 902 \\
  Random policy & $266.2 \pm 164.6$ & 240 & 2,290 & 499 \\
  \bottomrule
\end{tabular}
\end{center}

\begin{itemize}
  \item \textbf{The trained agent scores 6.7$\times$ the random policy}, and the two histograms barely overlap.
  \item \textbf{Selected vs.\ final network:} on validation, episode 4,000 is marginally ahead (1,857 vs.\ 1,835), on the test games the final network has a slightly higher mean (+52, about 1.7 standard errors) but more than twice the spread. 30 validation games cannot separate them. I report the selected network, as fixed before looking at the test games.
  \item \textbf{After episode 4,000, the validation score plateaus} at 1,600--1,850 (Figure~\ref{fig:b-snap}): more episodes at $\epsilon = 0.01$ do not help this configuration any more.
  \item \textbf{For orientation:} the original DQN paper (Mnih et al., 2015) reports 2,311 for Ms.\ Pac-Man after 50M agent steps, about 11$\times$ the training used here, with a slightly different evaluation protocol (no sticky actions).
\end{itemize}

\subsection{Non-Standard Plot: Where Does Ms.\ Pac-Man Go?}

Ms.\ Pac-Man's position is read from her colour in the raw frame. The heatmaps count how many steps she spends at each position over 10 greedy games, drawn over the maze.

\fig{1.0}{b_heatmap.png}{Ms.\ Pac-Man's position over 10 greedy games (log scale) for the random policy and four snapshots.}{fig:b-heat}

\begin{itemize}
  \item \textbf{Random policy:} she stays around the start position in the centre and dies there.
  \item \textbf{Episode 500:} long routes, but only through the lower half of the maze.
  \item \textbf{Episode 2,500:} most of the maze, but the top corridors are still almost never reached.
  \item \textbf{Episodes 4,000 and 6,000:} the \textbf{whole maze}, including the top corridors and the tunnels at the left and right edge. The agent has learned to clear the board, not only to collect the dots near the start.
\end{itemize}

The second plot follows the Q-value through one greedy game of the selected network:

\fig{1.0}{b_q_trace.png}{$\max_a Q(s, a)$ during one greedy test game (score 1,660), with power pellets, ghosts or fruit eaten, and lost lives.}{fig:b-trace}

\begin{itemize}
  \item \textbf{When play starts, Q jumps to its maximum ($\approx$ 45)} and then \textbf{falls steadily as the dots are eaten}: fewer dots left means less reward to come. This is exactly what a Q-value should do.
  \item \textbf{A power pellet lifts Q} (step $\approx$ 350): eatable ghosts are worth extra rewards.
  \item \textbf{Before every lost life, Q drops to $\approx$ 0--2}: the network sees death coming. After the respawn, Q is flat at $\approx$ 10 during the pause, lower than at the start because fewer dots are left.
  \item \textbf{Unexpected: during the first 67 steps (the start delay, in which the game ignores all actions), Q is $\approx -11$}, although all rewards are $\geq 0$. This is examined in Test~6.
\end{itemize}

\subsection{Edge Cases}

\begin{center}
\small
\begin{tabularx}{\textwidth}{@{}p{3.6cm}L@{}}
  \toprule
  \textbf{Edge case} & \textbf{Handling and check} \\
  \midrule
  Sticky actions & Kept, as v5 prescribes. Even the greedy policy is therefore stochastic. Without them, every greedy game is identical (Test~9). \\
  Flickering sprites & The ghosts are not drawn in every frame; with 4 stacked frames each ghost is visible in at least one (Figure~\ref{fig:b-pre}). \\
  Life loss vs.\ game over & A lost life is terminal for the TD target, but the game is not reset; only the third life ends the episode (Test~3 \pass). \\
  Time limit & A truncated game is not terminal. Its last step has no stored next state, so it is never sampled instead of being treated as terminal (Test~2 \pass). \\
  Frame stack at a game start & Filled with the first frame. In the buffer, stacks never reach back into the previous game but repeat the first frame, exactly like the online stack (Test~2 \pass). \\
  Buffer wrap-around & A sample needs all 4 frames of its stack still stored and its next state already stored (Test~2 \pass). \\
  Reward scale & Clipped to $\{0, 1\}$ for learning, raw score reported. Q is bounded by 100 (Test~6a \pass). \\
  int8 arithmetic of the preprocessing & \texttt{img.sum} gives unsigned integers, so \texttt{img // 3 - 128} wraps around; the cast to int8 still gives the right value (Test~1 \pass). \\
  Overestimation & Double DQN, slow target network, targets without gradients (Test~6). \\
  Crash during the overnight run & Atomic checkpoints with the replay buffer and all generator states; per-episode seeds; deterministic GPU algorithms (Test~10 \pass). \\
  Model selection on the test set & Snapshots are compared on validation seeds (10,000+), the test uses seeds 0--499. \\
  Memory & 1M frames as int8: 7\,GB of RAM; the full buffer is loaded from a checkpoint without a second copy. \\
  \bottomrule
\end{tabularx}
\end{center}

\subsection{Tests for Validating Correctness}

All tests use the selected network unless stated otherwise. Tests 1--3 check that the agent learns from the right data (the inputs, transitions and terminal signals the algorithm depends on), Tests 4--10 check the learning behaviour itself. Each criterion was fixed before the run. \textbf{23 of 24 checks pass.}

\begin{center}
\small
\begin{tabularx}{\textwidth}{@{}p{2.5cm}Lp{3.6cm}@{}}
  \toprule
  \textbf{Test} & \textbf{What is checked, and why it validates the algorithm} & \textbf{Result} \\
  \midrule
  1. Preprocessing & Frame is 88$\times$80 int8, equals a signed reference computation, Ms.\ Pac-Man's pixels become $-128$. \emph{Why:} the int8 cast wraps around; a wrong sign would hide Ms.\ Pac-Man from the network. & 3/3 \pass \\
  2. Replay buffer & Small buffer (64 slots, 421 frames, wrapped around): sampled stacks of $s$ and $s'$ equal the online stacks, no stack mixes two games, truncated last steps are never sampled. \emph{Why:} index-based stacking is the most error-prone part; a mismatch would train on states the agent never sees. & 5/5 \pass \\
  3. Life loss & A random game: 3 TD-terminal steps, only the last one ends the game, no reset in between. \emph{Why:} the terminal signal decides where the TD target stops bootstrapping. & 3/3 \pass \\
  4. Learning & Last 500 training episodes $> 2\times$ the first 500; test score $> 2\times$ random. \emph{Why:} the grading criterion; improvement must come from learning, not from chance. & 296 $\to$ 1,848; 1,791 vs.\ 266 \pass \\
  5. Progress & Validation score of the 13 snapshots rises with training (Spearman $\rho$); best snapshot $> 3\times$ the untrained one. \emph{Why:} greedy play on fixed games, independent of $\epsilon$ and of the training noise. & $\rho = 0.93$; 70 $\to$ 1,857 \pass \\
  6. Q vs.\ actual return & (a) max Q $< 1/(1-\gamma) = 100$; (b) correlation of Q with the actual discounted return $G$ along greedy games $> 0.5$; (c) Q in the 10 steps before a lost life $< 80\,\%$ of its mean. \emph{Why:} Q must estimate the return, not only rank actions. & (a) 57.3 \pass\ (b) 0.46 \fail\ (c) 4.5 vs.\ 12.5 \pass \\
  7. $\gamma = 0$ run & Changed parameter: with $\gamma = 0$, Q must equal the probability of a reward, so all Q in $[-0.1, 1.1]$ and mean Q = reward rate. \emph{Why:} tests the target without bootstrapping, where the right answer is known. & $[-0.03, 0.34]$; 0.053 vs.\ 0.052 \pass \\
  8. First 10 steps & Q of all actions, action, reward and TD target step by step; mean $|\text{TD error}| < 1$ over a full game. \emph{Why:} a trained network must be close to a fixed point of its own Bellman update. & 0.48 \pass \\
  9. Harder evaluation & Reduced and added stochasticity: 100 games without sticky actions (deterministic) and with $\epsilon = 0.05$, both $> 2\times$ random. \emph{Why:} the policy must not depend on one lucky sequence of events. & 1,710 $\pm$ 0; 1,782 $\pm$ 420 \pass \\
  10. Exact resume & 6 episodes in one go vs.\ 3 episodes, ``crash'', resume, 3 more: identical log and weights. \emph{Why:} the reported run was resumable; a resume must not change the result. & 2/2 \pass \\
  \bottomrule
\end{tabularx}
\end{center}

\fig{1.0}{b_test_q_calibration.png}{Test 6: predicted $Q(s_t, a_t)$ vs.\ the actual discounted return $G_t$ of the clipped rewards (cut at every lost life) over 10 greedy test games (left), and along one game (right).}{fig:b-calib}

\begin{itemize}
  \item \textbf{Test 6 (Q vs.\ return):} unlike part (a), these roll-outs are greedy, so Q and $G$ refer to the same policy. During play, Q follows $G$ closely, but the start of play is clearly \textbf{overestimated} (44.8 vs.\ $\approx$ 27 in all 10 games). Over all steps outside the start delay, the mean Q is 14.4 vs.\ a mean $G$ of 9.7.
  \item \textbf{Test 7} checks the target computation in isolation from bootstrapping: Q matches the reward rate to 0.001.
  \item \textbf{Test 8:} the literal first 10 steps fall into the start delay: all 9 Q-values lie within 0.02 of each other ($\approx -11.5$), which is correct in so far as the action does not matter there. The table therefore also shows the first 10 steps after the delay, when Ms.\ Pac-Man starts to eat (test game 0; $y = r + \gamma \max_a Q(s', a)$ with the same network):
\end{itemize}

\begin{center}
\footnotesize
\setlength{\tabcolsep}{4.5pt}
\begin{tabular}{@{}l*{10}{r}@{}}
  \toprule
  Step $t$ & 67 & 68 & 69 & 70 & 71 & 72 & 73 & 74 & 75 & 76 \\
  \midrule
  Action & \action{UP} & \action{UP} & \action{UP} & \action{DOWN} & \action{DOWN} & \action{DOWN} & \action{UPRIGHT} & \action{DOWN} & \action{NOOP} & \action{NOOP} \\
  $\max_a Q(s_t, a)$ & 44.82 & 44.04 & 44.08 & 44.12 & 43.12 & 43.10 & 43.49 & 42.59 & 42.69 & 43.05 \\
  $Q(s_t, \text{\action{LEFT}})$ & 44.58 & 43.72 & 43.30 & 43.13 & 42.14 & 42.82 & 43.39 & 42.54 & 42.59 & 43.05 \\
  Raw reward & 10 & 0 & 0 & 10 & 0 & 0 & 10 & 0 & 0 & 0 \\
  Clipped reward $r$ & 1 & 0 & 0 & 1 & 0 & 0 & 1 & 0 & 0 & 0 \\
  TD target $y$ & 44.60 & 43.64 & 43.68 & 43.69 & 42.67 & 43.06 & 43.17 & 42.26 & 42.62 & 42.77 \\
  TD error $y - Q(s_t, a_t)$ & $-0.22$ & $-0.40$ & $-0.40$ & $-0.44$ & $-0.45$ & $-0.04$ & $-0.33$ & $-0.33$ & $-0.07$ & $-0.29$ \\
  \bottomrule
\end{tabular}
\end{center}

\begin{itemize}
  \item \textbf{The Bellman update is visible step by step.} Whenever a dot is eaten ($r = 1$ at steps 67, 70, 73), Q drops by about 1 in the next step (44.8 $\to$ 44.0, 44.1 $\to$ 43.1, 43.5 $\to$ 42.6): one reward fewer is left to collect. Between two dots, Q rises again slightly as the next dot comes closer and is discounted less.
  \item \textbf{The TD errors are small but all negative} ($-0.04$ to $-0.45$): the network predicts slightly more than its own next-step target, the same overestimation at the start of play that Test~6 found. Over the whole game, the mean $|\text{TD error}|$ is 0.48 \pass.
  \item \textbf{The policy decides on tiny margins.} At step 67, \action{LEFT} is only 0.24 below the best action; the median action gap over the game is 0.028. As in part (a), most decisions are close calls, which is why the dueling head (learning $V$ separately) suits the game.
  \item \textbf{Test 9:} with 5\,\% random actions, the mean barely changes (1,782). \textbf{Without sticky actions, all 100 games are identical (1,710 $\pm$ 0):} the emulator is otherwise deterministic, so the whole spread of the evaluation comes from the sticky actions.
\end{itemize}

\begin{notebox}
\textbf{Why Test 6b fails: the start delay.} At the start of every game, Ms.\ Pac-Man cannot move for 67 agent steps. In these steps, the network predicts $Q \approx -11$, although with rewards $\geq 0$ every true Q-value is positive ($\approx 0.99^{67-t} \cdot 27$). These 7.6\,\% of the steps form the horizontal band at $Q \approx -11$ in Figure~\ref{fig:b-calib}. \textbf{Without them, the correlation is 0.80}, well above the threshold.

\smallskip
The delay states are not frozen frames: the ghosts move in a cycle of 16 steps, so the same stack appears four times at different distances to the start of play. The negative value appears in every snapshot and moves with training ($-0.6$ at episode 500, $-11.5$ at episode 4,000, $-2.0$ at episode 6,000), and the TD errors there are consistently positive ($+0.1$), so the update pulls it in the right direction but does not win against the generalisation from the rest of the network. I could not explain the cause with the remaining time. \textbf{It does not affect the score}, since the game ignores every action during the delay, but it is a real error of the value function.
\end{notebox}

\FloatBarrier
\begin{summarybox}{Summary of Part (b)}
\begin{itemize}
  \item \textbf{Model:} Double DQN with a dueling head, 4 stacked frames, 1M replay buffer, slow Polyak target network, Huber loss with clipping and clipped rewards; 6,000 episodes (4.46M agent steps, 5.8 hours).
  \item \textbf{Result:} the selected network (episode 4,000) scores \textbf{1,790.6 $\pm$ 286.8} over 500 greedy test games, 6.7$\times$ the random policy (266.2). It clears the whole maze instead of staying near the start, and its Q-value reacts sensibly to dots, power pellets and approaching death.
  \item \textbf{Validation:} 23 of 24 checks pass, including a bit-identical resume after a simulated crash. The failed check reveals a negative Q-value during the start delay, which does not affect play. Apart from it, Q follows the actual returns ($r = 0.80$) but overestimates them, especially at the start of play.
\end{itemize}
\end{summarybox}
